\section{Silhouette initialization: formulation and algorithm}
\label{sec:silhouette}
\subsection{Purpose and information available to the estimator}
The comparator supplies a large-displacement starting transformation from the
\emph{observable extent of the specimen}. It does not match contour landmarks,
describe the speckle, or estimate a nonlinear interior displacement from the
boundary. Its four unknowns are two axis-aligned scale factors and two
translations. This deliberately restricted hypothesis tests whether the local
solver can recover useful states once the order of magnitude of the motion is
provided by the images.

Only the reference image $I_0$ and the current target $I_1$ enter this
initialization. Neither the target index, a prescribed strain, the displacement
files, nor a preceding target solution is used. The complete reference
preparation may be cached, but the target states are reset for every pair.
The final field is a set of independently resolved local affine states, not
the global diagonal map imposed as a final measurement.

\paragraph{Notation and reading guide.}
$I_j$ denotes image $j$, with $j=0$ for the reference and $j=1$ for the target;
$\Omega_j$ is its pixel domain. Vectors $x$ and $q$ are positions in the reference
and target, respectively. $I$ without an image subscript denotes the $2\times2$
identity matrix. $d$ is a displacement in pixels; $G=\nabla d$ is its dimensionless
$2\times2$ gradient. The mask, box, support and propagation figures below follow
these quantities through the executed algorithm.

\subsection{Image-derived material masks}
Pixel centres have integer coordinates $(x,y)$, with $x$ increasing to the right
and $y$ downwards. In each image separately, the background value is read from
the bottom-right pixel, $b_j=I_j(W_j-1,H_j-1)$. Let
$E_j=\{(x,y):I_j(x,y)=b_j\}$. We label its four-connected components
$C_{jk}$ and define
\begin{equation}
 M_j=\Omega_j\setminus\bigcup_{k:\,|C_{jk}|\ge128}C_{jk}.
 \label{eq:silhouette-mask}
\end{equation}
Thus large constant-background regions, including sufficiently large internal
holes, are excluded. Small islands with the same quantized intensity are kept:
otherwise individual speckle pixels would perforate the mask and remove many
usable supports. The threshold is a component area in native pixels, not an
intensity tolerance. There is no morphological closing, learned segmentation,
subpixel contour fit or analytical material mask in this operation.

The reference mask $M_0$ defines the region of interest (ROI), meaning the
domain eligible for correlation. A window is eligible only when its full
reference support is contained in that mask. The target mask $M_1$ supplies
its bounding box; it is not passed as a target-material validity constraint
to the solver. Target-image bounds are checked by the solver, but material
occlusion is not independently certified. Exact equality to a single background
intensity is suitable for these rendered images, not a general segmentation
rule for noisy experimental backgrounds.

Figure~\ref{fig:mask-steps} separates the observed image, the binary material
mask and the eligible supports. If $\Xi_i$ is the set of pixel offsets in window
$i$, full-support eligibility is the Boolean product
\begin{equation}
 V_i=\prod_{\xi\in\Xi_i}{\bf1}_{M_0}(x_i+\xi).
 \label{eq:full-support}
\end{equation}
Here ${\bf1}_{M_0}$ is one inside the reference material mask and zero outside.
This is a support test, not an estimate of target visibility or matching accuracy.
\begin{figure}[H]\centering
\includegraphics[width=\linewidth]{silhouette_mask_steps.pdf}
\caption{From reference pixels to admissible windows on the tested perforated
plate. White in the mask denotes retained material; green outlines show the
36 admissible 32-pixel supports. Holes are excluded from supports, not paired
as corresponding landmarks. All panels use native reference coordinates.}
\label{fig:mask-steps}\end{figure}

\subsection{Pixel-edge boxes and the four-parameter proposal}
For a nonempty mask, define the lower and upper edges of its bounding box:
\begin{align}
 \ell_j&=(\min_{M_j}x-\tfrac12,\ \min_{M_j}y-\tfrac12)^\top,\\
 h_j&=(\max_{M_j}x+\tfrac12,\ \max_{M_j}y+\tfrac12)^\top,
 \qquad L_j=h_j-\ell_j.
\end{align}
A run of $n$ pixels therefore has extent $n$, not $n-1$. This half-pixel
convention is important when the scale is large. Matching the two pairs of
axis-aligned box edges gives
\begin{equation}
 S=\begin{bmatrix}L_{1x}/L_{0x}&0\\0&L_{1y}/L_{0y}\end{bmatrix},
 \qquad t=\ell_1-S\ell_0,\qquad \varphi_0(x)=Sx+t.
 \label{eq:silhouette-map}
\end{equation}
Indeed $\varphi_0(\ell_0)=\ell_1$ and $\varphi_0(h_0)=h_1$.
These conditions identify a map between boxes; they do not establish material
correspondence between their extremal pixels. The initial displacement and
its gradient are
\begin{equation}
 u_0(x)=(S-I)x+t,\qquad G_0=\nabla u_0=S-I.
\end{equation}
No rotation or shear is estimated at this stage. Internal holes affect window
eligibility but their relative positions are not fitted as landmarks.

For example, the strip reference has vertical edges $[-0.5,511.5]$ and target
500 has edges $[-0.5,5631.5]$. Consequently $s_y=11$, $t_y=5$ pixels and
$v_0(y)=10y+5$. Using $v_0=10y$ would introduce a five-pixel offset despite
matching an apparently fixed top edge. For plate target 191, the edges change
from $[3.5,214.5]$ to $[3.5,1018.5]$, yielding
$s_y=1015/211\simeq4.81043$ and $t_y\simeq-13.33649$ pixels.
Both examples have $s_x=1$ and $t_x=0$. These are observed envelope ratios,
not known local stretches.

The geometry is particularly clear when written relative to the lower edge:
\begin{equation}
 \boxed{\varphi_0(x)=\ell_1+S(x-\ell_0).}
 \label{eq:relative-box}
\end{equation}
Subtracting $\ell_0$ changes origin, multiplying by $S$ changes the two extents,
and adding $\ell_1$ places the result in the target image. For either coordinate
$\alpha\in\{x,y\}$, this is exactly the solution of the two edge constraints
\begin{equation}
 \begin{bmatrix}\ell_{0\alpha}&1\\h_{0\alpha}&1\end{bmatrix}
 \begin{bmatrix}s_\alpha\\t_\alpha\end{bmatrix}
 =\begin{bmatrix}\ell_{1\alpha}\\h_{1\alpha}\end{bmatrix}.
\end{equation}
There is no interior photometric minimization in this step.
\begin{figure}[H]\centering
\includegraphics[width=.92\linewidth]{silhouette_box_geometry.pdf}
\caption{Construction of the vertical map for plate target 191. The values are
measured pixel-edge coordinates from the saved pair. Each panel has its own
vertical scale for readability; the annotated lengths, not the drawn heights,
define the ratio. The lower edge is fixed here at 3.5 pixels, which still gives
a nonzero translation coefficient when the map is expressed about the image origin.}
\label{fig:box-geometry}\end{figure}

\subsection{Conversion to the six local solver parameters}
For a window centred at $x_i=(x_i,y_i)^\top$, the inverse-compositional
Gauss--Newton (IC-GN) solver uses
\begin{equation}
 p_i=(u_i,u_{x,i},u_{y,i},v_i,v_{x,i},v_{y,i})^\top.
\end{equation}
Its local forward warp, for a reference offset $\xi$, is
$q=x_i+(u_i,v_i)^\top+(I+G_i)\xi$, where
$G_i=\left[\begin{smallmatrix}u_{x,i}&u_{y,i}\\v_{x,i}&v_{y,i}\end{smallmatrix}\right]$.
The box proposal supplies
\begin{equation}
 p_i^0=\big((s_x-1)x_i+t_x,\ s_x-1,\ 0,\
                  (s_y-1)y_i+t_y,\ 0,\ s_y-1\big)^\top.
 \label{eq:silhouette-six}
\end{equation}
All six components are passed to the solver, including the nonzero diagonal
gradients. The native target pixels are sampled by the local warp: the full
image is not first rescaled and then treated as an undeformed translation
problem. During the solve, all six parameters may change, including the two
cross derivatives initialized to zero. This is the mechanism by which local
heterogeneity can be recovered from a simple global proposal.

Figure~\ref{fig:local-support} illustrates the distinction between an initial
support and a solved support. At initialization the following identity holds:
\begin{equation}
 \underbrace{x_i+d_i^0+(I+G_i^0)\xi}_{\text{local affine sampling}}
 =\underbrace{S(x_i+\xi)+t}_{\text{global box proposal}},
 \qquad d_i^0=(S-I)x_i+t,\quad G_i^0=S-I.
 \label{eq:local-global-identity}
\end{equation}
After a local solve there is no requirement that $G_i=G_i^0$ or that all windows
share the same $G_i$. The pulled-back target patch is
\begin{equation}
 P_{1,i}(\xi;p_i)=I_1\big(x_i+d_i+(I+G_i)\xi\big),\qquad
 P_{0,i}(\xi)=I_0(x_i+\xi).
 \label{eq:patch-pullback}
\end{equation}
For nonconstant sampled patches, their zero-mean normalized cross-correlation is
\begin{equation}
 \rho_i(p_i)=
 \frac{\sum_{\xi\in\Xi_i}(P_{0,i}-\overline P_{0,i})(P_{1,i}-\overline P_{1,i})}
 {\sqrt{\sum_{\xi\in\Xi_i}(P_{0,i}-\overline P_{0,i})^2}
  \sqrt{\sum_{\xi\in\Xi_i}(P_{1,i}-\overline P_{1,i})^2}}.
 \label{eq:patch-correlation}
\end{equation}
Bars denote means over the support. This equation defines the photometric
correlation diagnostic and its 0.65 threshold; correlation alone is not the
complete numerical acceptance rule, nor a ground-truth correctness criterion.
A translation-only initialization would replace $(I+G_i^0)\xi$ by $\xi$ and
sample the wrong material support under strong stretching.
\begin{figure}[H]\centering
\includegraphics[width=\linewidth]{silhouette_local_support.pdf}
\caption{One saved accepted window from plate target 191. Blue marks its reference
support. Orange is its image under the box proposal; green is its image under
the final solved local affine state. Centres and outlines come from the archived
parameters, not an analytical deformation. The selected window is the accepted
centre nearest the centroid of eligible centres, using a deterministic geometric rule.}
\label{fig:local-support}\end{figure}

\subsection{Status-based selection and confirmation of seeds}
Two hypotheses are tested in the fixed order: zero displacement/gradient,
then the silhouette proposal. On the regular eligible-window grid, every
second row and every second column are sampled. The identical sparse subset
is solved under each hypothesis. Finite solver-accepted states are counted;
the hypothesis with the larger count wins. A tie selects the first, zero
hypothesis. No displacement truth, correlation-error ground truth or target
index enters this ranking. The accepted counts in this selection are not
yet the final propagated counts.

Writing $\mathcal S$ for the common sparse set and $a_i(h)$ for the Boolean
finite, solver-accepted status after solving hypothesis $h$, the selection is
\begin{equation}
 h^*=\operatorname*{arg\,max}_{h\in\{\mathrm{zero},\mathrm{box}\}}
       \sum_{i\in\mathcal S} a_i(h),
 \label{eq:hypothesis-selection}
\end{equation}
with ties resolved by the displayed order. The status is measured after solving,
not by comparing the uncorrected initial patches alone.

Let $d_i=(u_i,v_i)^\top$ and let $j$ be an eight-neighbour of $i$. The first-order
prediction transported from $i$ to $j$ is
\begin{equation}
 \widehat d_{j\mid i}=d_i+G_i(x_j-x_i),\qquad
 \widehat G_{j\mid i}=G_i.
 \label{eq:silhouette-transport}
\end{equation}
Neighbour agreement uses the Euclidean displacement norm and the Frobenius
gradient norm:
\begin{equation}
 \|d_j-\widehat d_{j\mid i}\|_2\le8\ \mathrm{px},\qquad
 \|G_j-G_i\|_F\le0.25.
 \label{eq:silhouette-consistency}
\end{equation}
The implementation has two confirmation routes. A candidate anchor and an
already verified adjacent state are confirmed if they satisfy these tests.
Otherwise eligible, not-yet-verified neighbours are solved from transported
states; at least two successful and consistent neighbour solves confirm the
parent and those neighbours. Solver status and finite parameters are also
required. Batched confirmation is enabled and used when more than eight
candidate anchors exist; results are replayed in the original confirmation
order.

The argument \texttt{minimum\_anchors=3} only produces a diagnostic indicating
whether each eligible grid component contains three confirmed anchors. In
this comparator it does \emph{not} reject a pair or disable propagation.
It must not be described as an enforced three-anchor acceptance gate.

\subsection{Frontier propagation and its precise stopping rule}
Each propagation pass starts from an immutable snapshot of accepted windows.
For an unresolved eligible window, the first two available accepted neighbours
are considered in the fixed order up, left, right, down, then the four
diagonals. The transported proposal is solved on the images. A finite,
solver-accepted result is retained only if it also satisfies
Equation~\eqref{eq:silhouette-consistency} relative to that parent prediction.
The second parent is attempted only when the first has not produced an
accepted result. New states become parents on the next pass; accepted states
remain unchanged. Propagation stops if no frontier exists or a complete pass
adds no accepted window.

This is the existing confirmed-propagation policy: it is not a median over
all neighbours, not inverse-distance averaging, and not the historical
three-attempt frontier policy. There is no fixed three-attempt lifetime cap
here. Neighbourhoods use the eligible-window grid, not an explicit
pixel-level material-connectivity test along every link. The reference
support mask prevents solving windows over holes but does not prove that
every diagonal link avoids a narrow fissure. Discontinuity preservation is
therefore not established by this experiment.

Let $A^{(k)}$ be the accepted set at the start of pass $k$, and $N(i)$ the
ordered eight-neighbour set. The frontier is
\begin{equation}
 \mathcal F^{(k)}=\{i:V_i=1,\ i\notin A^{(k)},\ N(i)\cap A^{(k)}\ne\emptyset\}.
\end{equation}
After a parent proposal has been solved, denote its numerical status by $s_i$
and its returned state by $(d_i,G_i)$. The propagation acceptance predicate is
\begin{equation}
 b_{i|j}=s_i\ \land\ \operatorname{finite}(p_i)\ \land\
 \big(\|d_i-\widehat d_{i|j}\|_2\le8\big)\ \land\
 \big(\|G_i-G_j\|_F\le0.25\big).
\end{equation}
A window is added if one of its attempted parents passes this predicate.
The actual archived generations and parent links are shown in
Figure~\ref{fig:propagation-steps}. They are not displacement vectors.
\begin{figure}[H]\centering
\includegraphics[width=\linewidth]{silhouette_propagation_steps.pdf}
\caption{Measured progression on plate target 191 with 32-pixel supports.
Confirmed seeds are generation zero. The middle panel retains states accepted
by generation two; the final panel shows recorded parent links and acceptance
generations. Orange crosses denote final rejected windows in all three panels.
Blank centres in an early panel can still be recovered later. No unrecorded
sparse-trial locations or synthetic displacement truth are reconstructed.}
\label{fig:propagation-steps}\end{figure}

\subsection{Frozen settings and observed stages}
The solver operates on native grayscale pixels with affine motion, bilinear
interpolation, at most 30 iterations, tolerance $10^{-3}$, and minimum
zero-mean normalized cross-correlation (ZNCC) 0.65. Its separate outlier policy
is disabled; the above neighbour consistency checks remain active. The
64-pixel experiment uses support and stride 64; the separate 32-pixel
experiment uses both 32. No angular bank, learned residual, separate local
correlation search, regional retry, or coarse-to-fine support transfer is
used by this comparator.

\begin{center}\small
\begin{tabular}{lrrrr}\toprule
Pair and support & Zero trial & Box trial & Confirmed seeds & Final accepted\\
\midrule
Plate 191, 32 pixels & 0 & 4 & 7 & 18/36\\
Strip 500, 64 pixels & 0 & 8 & 35 & 109/128\\
\bottomrule\end{tabular}
\end{center}
Both examples require six propagation passes, including the stopping pass
when applicable. Confirmed seeds can outnumber sparse trial acceptances
because neighbour confirmation performs additional solves. These counts
come directly from the saved per-pair records. They show that the final
coverage is not simply the fraction accepted under an uncorrected box map.
They do not establish the correctness of the accepted displacement fields.

\subsection{Interpretation, failure modes and computational scope}
The useful role of the silhouette is to turn a difficult global texture
search into a plausible starting scale and translation. If several windows
fall inside the solver's convergence domain, their solved gradients can
transport initializations locally, without forcing the entire specimen to
follow the diagonal box map. Conversely, a correct envelope does not guarantee
any correct interior seed. Different non-affine fields can share the same
bounding box, including fields with severe interior localization.

Cropping, occlusion, detached components, rotation, shear, an incorrectly
identified corner background, or a few remote segmentation outliers can bias
the extrema. The method has no built-in uncertainty estimate for these
effects. For one coordinate $s=L_1/L_0$, first-order perturbation gives
$\delta s\simeq\delta L_1/L_0-L_1\delta L_0/L_0^2$;
$\delta t=\delta\ell_1-s\delta\ell_0-\ell_0\delta s$.
Short reference extents and large scale ratios amplify boundary uncertainty.
These limitations preclude presenting the method as a general replacement
for feature matching or as a deformation measurement from contours.

Mask labelling and box extraction are host-side array operations. Hypothesis
selection, confirmation and propagation are orchestrated in Python; the
local solves use the WebGPU graphics interface and WebGPU Shading Language
(WGSL) kernels. Parameter/status arrays cross the host--device boundary
between calls. This experiment is therefore not a wholly graphics-resident
pipeline, and it does not establish real-time performance. Its reproducible
role is a non-learned, image-only control within the same local solver.

The new evidence motivates adding this proposal as an optional candidate
when complete specimen silhouettes are available. It does not establish
novelty or general superiority. Independent displacement truth, segmentation
perturbation tests, and paired ablations of proposal-only, confirmation and
propagation are still required before attributing a quantified accuracy gain
to the silhouette itself.
